\section{Results}
\label{sec:results}

Unless stated otherwise, Stage~I values are descriptive means across frozen model or model--regime combinations; model families were not treated as independent inferential replicates. Audit contrasts were paired on the same endpoint, split, seed, model, and regime, and seed~99 was excluded from scientific conclusions. Stage~II development and final results are reported separately so that architecture-selection evidence is not confused with independent validation.

\subsection{Chemical shift changed point performance and calibration}
\label{sec:results_performance}

Random splitting generally produced the strongest point performance, but the size and form of degradation depended on the endpoint (Figure~\ref{fig:shift_performance}). BBBP mean ROC--AUC declined from 0.886 under random splitting to 0.851 under scaffold and 0.864 under similarity-cluster splitting. Balanced accuracy decreased from 0.789 to 0.736 and 0.748, respectively. ClinTox was more difficult: ROC--AUC ranged from 0.675 to 0.703 and balanced accuracy from 0.537 to 0.575, with the random split again strongest.

The two regression endpoints separated a split effect from a weak-model effect. ESOL mean $R^2$ was 0.488 under random splitting and 0.451 under cluster splitting but fell to 0.089 under scaffold separation; RMSE concurrently increased from 1.441 to 2.035. Lipophilicity mean $R^2$ was already negative under random splitting ($-0.175$) and deteriorated under scaffold ($-0.503$) and cluster ($-0.617$) designs. Conformal coverage close to 0.90 for this endpoint therefore could not be read as evidence of useful point prediction.

Positive-probability calibration also favored random splitting within endpoint. BBBP expected calibration error increased from 0.075 under random splitting to 0.096 under scaffold splitting; the corresponding ClinTox values were 0.051 and 0.057. Brier score and negative log-likelihood showed the same within-endpoint direction. Their smaller raw ClinTox values were not interpreted across endpoints because both scores depend on prevalence. Overall, the structured splits did not produce one uniform penalty: they changed discrimination, explained variance, and calibration in endpoint-specific ways.

\begin{figure}[!htbp]
\centering
\includegraphics[width=\textwidth]{figure_2_shift_performance.pdf}
\caption{\textbf{Point prediction and calibration under chemical distribution shift.} Descriptive means across frozen model or model--regime combinations are shown for random, label-blind scaffold, and label-blind similarity-cluster splits. (A) ROC--AUC for BBBP and ClinTox. (B) Balanced accuracy. (C) $R^2$ for ESOL and Lipophilicity. (D) Positive-probability expected calibration error. Dashed lines denote chance discrimination or zero explained variance where applicable. Model-specific cross-seed uncertainty is reported in the Supporting Information.}
\label{fig:shift_performance}
\end{figure}

\subsection{Aggregate conformal coverage concealed class failure and low information}
\label{sec:results_coverage}

At nominal 90\% coverage, overall classification coverage was close to target while class-specific behavior diverged sharply (Figure~\ref{fig:reliability_illusion}A,B). For BBBP, marginal overall coverage was 0.904--0.907 across splits, but positive coverage was 0.971--0.978 and negative coverage only 0.664--0.694. Density-ratio weighting changed these values only modestly. Mondrian calibration reduced the gap, although mean set size increased by approximately 0.208--0.232 and ambiguity increased by approximately 0.207--0.232.

ClinTox exposed the more consequential failure. Marginal overall coverage was 0.897--0.902, yet positive-class coverage was 0.068 under scaffold, 0.069 under cluster, and 0.099 under random splitting; negative-class coverage was 0.954--0.968. Shift weighting left overall coverage at 0.890--0.912 while positive coverage remained only 0.073--0.158. These sets looked precise rather than conservatively broad: mean size was below one because approximately 3.5--5.0\% were empty, and two-label sets were rare.

Mondrian calibration raised ClinTox positive coverage by 0.841--0.872 relative to marginal prediction, with all eight frozen model--regime combinations showing model-specific cross-seed intervals entirely above zero. Positive coverage reached approximately 0.926--0.940 and negative coverage approximately 0.895--0.900. The repair was informative about uncertainty but inefficient as a binary decision: 70.9--71.9\% of test compounds received $\{0,1\}$, and mean set size reached 1.71--1.72. Thus overall coverage alone obscured two opposite defects---rare-class undercoverage for marginal sets and low decisiveness for class-conditional sets.

The same distinction reappeared during public Tox21 development (Figure~\ref{fig:reliability_illusion}C). ECFP density weighting produced minimum-class coverage of 1.00 across the 30 primary endpoint--seed cells, but MacroCSY was 0 and ambiguity was 1.00: every output was a two-label set. The physicochemical, score-view, ordinary, and TAME summaries occupied a lower-coverage but substantially more informative region. This was the empirical reason that a transport view could not be admitted merely because it increased coverage.

\begin{figure}[!htbp]
\centering
\includegraphics[width=\textwidth]{figure_3_reliability_illusion.pdf}
\caption{\textbf{Aggregate validity can hide either class failure or uninformative prediction sets.} (A) ClinTox overall, positive-class, and negative-class coverage for marginal, ECFP shift-weighted, and Mondrian conformal prediction, averaged across the frozen audit designs. (B) Minimum class coverage versus ambiguous two-label rate across BBBP and ClinTox split designs; marker shape identifies the conformal method and color identifies the endpoint. The arrow indicates the direction of jointly greater coverage and lower ambiguity. (C) Public Tox21 development summaries across six primary endpoints and five seeds. Bubble area is proportional to ambiguity; the ECFP-weighted comparator attained complete coverage by returning two labels for every compound.}
\label{fig:reliability_illusion}
\end{figure}

\subsection{Regression conformal variants did not show a universal advantage}
\label{sec:results_regression}

Marginal, shift-weighted, and adaptive normalized regression intervals all produced coverage near 0.90, but neither added method consistently improved coverage and efficiency together. ESOL marginal coverage was 0.901 under random, 0.891 under scaffold, and 0.911 under cluster splitting. Paired shift-weighted and adaptive changes relative to marginal prediction were small ($-0.010$ to $+0.006$), and no comparison was supported across all four regression models. For Lipophilicity, shift weighting increased descriptive coverage by 0.006 under scaffold and 0.011 under cluster splitting, while increasing mean width by 0.174 and 0.234.

Adaptive normalization produced only modest interval heterogeneity: width coefficients of variation were 0.007--0.013 for ESOL and approximately 0.032--0.033 for Lipophilicity. Mean Spearman associations between interval width and realized error were weak, approximately 0.05--0.20. The audit therefore supplied no basis for assuming that estimated shift correction or local scaling would improve reliability across endpoints.

\subsection{Chemical similarity and selective rejection were endpoint dependent}
\label{sec:results_domain}

Maximum Tanimoto similarity to the training set ranked risk for some endpoints but not others (Figure~\ref{fig:domain_retention}A,B). BBBP similarity--risk correlations ranged from $-0.131$ to $-0.159$ and similarity--miscoverage correlations from $-0.120$ to $-0.127$. Lipophilicity showed the strongest associations: approximately $-0.22$ to $-0.33$ for risk and $-0.19$ to $-0.24$ for miscoverage. ClinTox correlations were weak, approximately $-0.02$ to $-0.06$. ESOL reversed under scaffold splitting, with correlations of approximately $+0.04$ for risk and $+0.08$ for miscoverage even though the random and cluster associations were negative. A fixed similarity threshold would therefore define different scientific contrasts across endpoint and split.

Selective prediction clarified why this mattered. At 50\% overall retention, confidence-based ClinTox rejection retained only 26.9--30.6\% of positive examples, whereas similarity-based rejection retained 44.0--46.1\% (Figure~\ref{fig:domain_retention}C). Confidence lowered ordinary-risk partial area by 0.021--0.028 but worsened balanced-risk partial area by 0.021--0.029. The apparent benefit partly reflected a change in who remained.

Similarity-based regression rejection was likewise not transferable as a universal rule (Figure~\ref{fig:domain_retention}D). It improved risk at 50\% retention for Lipophilicity by approximately 0.245--0.321 across all splits. For ESOL, it helped under random and cluster splitting but was worse than matched random rejection under scaffold splitting by approximately 0.101. Together, these results made two requirements explicit for Stage~II: transport evidence had to be view- and batch-specific, and a method could not earn confidence by deleting baseline labels.

\begin{figure}[!htbp]
\centering
\includegraphics[width=\textwidth]{figure_4_domain_retention.pdf}
\caption{\textbf{Chemical-domain and selective-retention diagnostics are endpoint dependent.} (A) Mean Spearman association between maximum training-set similarity and predictive risk. (B) Association between similarity and marginal conformal miscoverage. Negative values indicate increasing risk or miscoverage as similarity decreases. (C) ClinTox positive-class retention when 50\% of all samples remain after confidence- or similarity-based rejection; the dashed line is composition-neutral 50\% retention. (D) Improvement in regression risk at 50\% retention for similarity rejection relative to matched random rejection. Threshold and rejection analyses are sensitivity diagnostics, not universal chemical-domain rules.}
\label{fig:domain_retention}
\end{figure}

\subsection{Audit failures became frozen TAME constraints}
\label{sec:results_development}

The audit did not identify a single scalar uncertainty score to optimize. Instead, it identified constraints: protect the weakest class explicitly; assess coverage jointly with correct-singleton yield; require batch-specific evidence before applying estimated transport weights; preserve every ordinary baseline label; expose an exact fallback; and separate architecture development from independent labels. Figure~\ref{fig:tame_mechanism} shows how these constraints were encoded in TAME.

Development results were retained as design evidence only. Across 30 public primary cells, the frozen candidate increased mean class-0 coverage by 5.4394 percentage points, left mean class-1 coverage unchanged, and changed MacroCSY by $-4.9169$ points. It produced zero baseline-inclusion violations, zero empty sets, and zero newly created singletons. The score view passed in 30/30 cells and the physicochemical view in 24/30, meeting the 80\% all-view activity gate; exact ordinary fallback handled the remaining six. The ECFP diagnostic failed in 30/30 cells because of effective-sample-size and/or clipping failures and was excluded from the envelope without relaxing its thresholds. All eight frozen promotion checks passed before any final EPA label was acquired or parsed.

\begin{figure}[!htbp]
\centering
\includegraphics[width=\textwidth]{figure_5_tame_mechanism.pdf}
\caption{\textbf{Mechanism and invariants of the transport-audited multi-view envelope.} (A) Four frozen ECFP classifiers feed a logit router and ordinary class-conditional Mondrian baseline. (B) Physicochemical and score views estimate separate label-free source-to-target weights. (C) Each view must pass the complete frozen transport certificate; the envelope requires an all-view quorum. (D) If both views pass, label~0 is added only by consensus while every baseline label is preserved. If either view fails, TAME returns the repaired ordinary baseline. The construction emits no empty set and cannot create a new confident singleton.}
\label{fig:tame_mechanism}
\end{figure}

\subsection{Sealed EPA validation improved minimum-class coverage within the efficiency bound}
\label{sec:results_final}

The independent run completed all 12 official endpoints across five fresh seeds (60/60 endpoint--seed cells); the primary analysis comprised the prespecified six endpoints and 30 cells. Two of 647 external structures failed the frozen standardizer and were retained as explicit full-set exclusions. Endpoint-specific standardized overlaps with the Tox21 training cohort were removed according to the lock. Predictions, the label-blind transport audit, and their hashes were written before the final label file was opened.

TAME increased the endpoint--seed-equal mean minimum-class coverage by 1.3649 percentage points relative to ordinary routed Mondrian prediction (Figure~\ref{fig:epa_validation}A,B). The deterministic 2000-draw hierarchical-bootstrap 95\% interval was $+0.5827$ to $+2.0051$ points. Five of six endpoint means were positive: NR-AhR $+2.1644$, NR-ER $+1.2362$, SR-ARE $+2.1461$, SR-MMP $+1.1489$, and SR-p53 $+1.4936$ points. SR-ATAD5 had zero minimum-class-coverage change even though its class-0 coverage increased by 1.4363 points; the unchanged class remained the minimum.

The primary mean MacroCSY change was $-1.6067$ percentage points, inside the frozen $-5$-point non-inferiority bound (Figure~\ref{fig:epa_validation}C). Because the envelope only adds a label, this efficiency cost is the expected counterpart of converting some singleton predictions into explicit deferrals. Structural checks confirmed the predeclared invariants: every TAME set contained its ordinary baseline set, no TAME set was empty, and no new singleton was created. Accordingly, wrong-singleton exposure did not increase on the same rows.

Both approved transport views passed in 60/60 final endpoint--seed cells, whereas the high-dimensional ECFP diagnostic failed in 60/60 because its frozen effective-sample-size, clipping, and balance certificate was not satisfied (Figure~\ref{fig:epa_validation}D). The diagnostic comparator nevertheless attained mean minimum-class coverage of 0.998 across all 12 endpoints while its mean ambiguity was 0.997 and MacroCSY only 0.002. This independent secondary result reproduced the central audit lesson: a coverage number can be nearly perfect because the output is nearly always uninformative.

The predeclared interpretation was therefore a coverage gain with efficiency non-inferiority in the locked Tox21 evaluation. It was not an automatic declaration of scientific superiority. The result does not establish conditional coverage, exact validity for estimated density ratios, transfer to arbitrary chemical shifts, or a general rule that label~0 should be protected in other assay families.

\begin{figure}[!htbp]
\centering
\includegraphics[width=\textwidth]{figure_6_epa_validation.pdf}
\caption{\textbf{Independent EPA validation under the prediction-to-label firewall.} (A) Mean TAME-minus-ordinary change in minimum-class coverage for each of the six prespecified primary endpoints across five fresh seeds. (B) Primary endpoint--seed-equal estimate and deterministic hierarchical-bootstrap 95\% interval (2000 draws; seed~44021). (C) Endpoint-level MacroCSY changes and their aggregate; the dashed line is the frozen $-5$-percentage-point non-inferiority bound. (D) Final activity of the two approved label-free transport views and the ECFP diagnostic across all 12 endpoints and five seeds. Predictions and transport audits were hashed before labels were opened.}
\label{fig:epa_validation}
\end{figure}
